\section*{Bab \ref{ch:g7:literal} --- Bentuk Aljabar}

\begin{solution}{exo:g7:literal:1}
$7x$; \quad $3y$; \quad $5ab$; \quad $x^2$; \quad $2(a + 4)$.
\end{solution}

\begin{solution}{exo:g7:literal:2}
$5ab = 5 \times a \times b$; \quad
$3(x+2) = 3 \times (x + 2)$; \quad
$x^2 + 4x = x \times x + 4 \times x$; \quad
$2a^3 = 2 \times a \times a \times a$.
\end{solution}

\begin{solution}{exo:g7:literal:3}
Untuk $x = 3$: \ $5 \times 3 + 1 = 16$; \quad $3 \times 3 = 9$; \quad
$2 \times 9 - 3 = 15$; \quad $10 - 6 = 4$.
\end{solution}

\begin{solution}{exo:g7:literal:4}
$3a + 2b = 12 + 3 = 15$; \qquad
$ab + 5 = 4 \times 1.5 + 5 = 6 + 5 = 11$.
\end{solution}

\begin{solution}{exo:g7:literal:5}
$8x + 5x = 13x$; \quad
$9a - 4a + a = 6a$; \quad
$6y + 2 + 3y + 7 = 9y + 9$; \quad
$5t + 3 - 2t - 3 = 3t$.
\end{solution}

\begin{solution}{exo:g7:literal:6}
Persegi: $P = 4c$. \omterm{def:g6:shapes:triangles}{Segitiga sama kaki}: $P = b + 2s$. Roti:
$p = 1.20\,n$ euro.
\end{solution}

\begin{solution}{exo:g7:literal:7}
$2n$: dua kali lipatnya; \ $n + 2$: dua lebih banyak daripada $n$;
\ $n^2$: kuadratnya; \ $\frac n2$: setengahnya.
\end{solution}

\begin{solution}{exo:g7:literal:8}
$3(x+4) = 3x + 12$: \emph{benar} untuk setiap $x$ (sifat distributif
--- ini sebuah bukti, tidak perlu diuji).

$2x + 5 = 7x$: ujilah $x = 2$: kiri $9$, kanan $14$ --- \emph{salah}
secara umum (ia hanya berlaku untuk $x = 1$).

$x + x = x^2$: ujilah $x = 3$: kiri $6$, kanan $9$ --- \emph{salah}
secara umum ($x + x = 2x$).
\end{solution}

\begin{solution}{exo:g7:literal:9}
\emph{1.} Panjangnya: $w + 3$. \omterm{def:g6:measure:perimeter}{Kelilingnya}:
\[
P = 2\bigl(w + (w + 3)\bigr) = 2(2w + 3) = 4w + 6 .
\]

\emph{2.} Untuk $w = 5$: secara langsung, sisinya $5$ dan $8$, jadi
$P = 2 \times 13 = 26$ cm; dengan rumusnya, $4 \times 5 + 6 = 26$ cm.
Keduanya cocok.
\end{solution}

\begin{solution}{exo:g7:literal:10}
Mengikuti langkahnya dengan huruf $n$:
\[
n \ \to\ n + 6 \ \to\ 3(n + 6) = 3n + 18 \ \to\ 3n + 18 - 18 = 3n
\ \to\ \frac{3n}{n} = 3 .
\]
Semua orang memperoleh $3$.
\end{solution}

\begin{solution}{exo:g7:literal:11}
Sebagai \omterm{ex:g1:subtraction:difference}{selisih}: $A = a^2 - b^2$. Dengan memotong huruf L itu menjadi
dua persegi panjang: satu jalur selebar penuh setinggi $a - b$ dan, di
atasnya, sebuah persegi panjang yang lebih sempit:
\[
A = a(a - b) + b(a - b)
\]
(persegi panjang bawah selebar $a$ dan setinggi $a - b$, persegi
panjang atas selebar $a - b$ dan setinggi $b$ memberi ragam
$a(a-b) + (a-b)b$ --- sama saja). Untuk $a = 5$, $b = 2$:
$a^2 - b^2 = 25 - 4 = 21$, dan
$5 \times 3 + 2 \times 3 = 15 + 6 = 21$. Kedua bentuknya cocok ---
dan memang selalu cocok: inilah kesamaan $a^2 - b^2 = (a + b)(a - b)$
pada \cref{ch:g9:algebra} yang sedang menyamar.
\end{solution}

\begin{solution}{pb:g7:literal:1}
\textbf{1.} Gambar $1$ sampai $4$ memerlukan $4$, $7$, $10$, $13$
batang.

\textbf{2.} Persegi yang baru berbagi satu sisi dengan gambar
sebelumnya, jadi hanya $3$ dari $4$ sisinya yang berupa batang baru.
Dari gambar $4$ ($13$ batang), enam persegi lagi memakan
$6 \times 3 = 18$: gambar $10$ memerlukan $13 + 18 = 31$ batang.

\textbf{3.} Membaca barisnya dari kiri: ia bermula dengan satu batang
tegak, lalu masing-masing dari $n$ persegi itu menyumbang sisi atas,
sisi bawah dan sisi kanan --- $3$ batang untuk tiap persegi.
Seluruhnya: $3n + 1$; ``$+1$'' itu adalah batang tegak yang paling
awal. Periksa: $n = 1, 2, 3, 4$ memberi $4, 7, 10, 13$.

\textbf{4.} $4 + 3(n - 1) = 4 + 3n - 3 = 3n + 1$: rumus Amir
menjabar menjadi bentuk sederhana yang persis sama. Sepakat sepenuhnya.

\textbf{5.} Lina: $2n + (n + 1) = 3n + 1$ --- sama lagi. Tiga cara
memandang, satu bentuk sederhana. Gambar $100$:
$3 \times 100 + 1 = 301$ batang.

\textbf{6.} Hitungannya: $3$, $5$, $7$ batang. Setiap segitiga baru
bersandar pada sisi yang sudah ada dan hanya menambah $2$ batang,
mulai dari $3$: gambar $n$ memerlukan $3 + 2(n - 1) = 2n + 1$ batang.

\textbf{7.} Bongkarlah $2n + 1 = 49$: membuang batang yang pertama
menyisakan $48$, lalu setiap segitiga bernilai $2$:
$48 \div 2 = 24$ segitiga.

\textbf{8.} Pada sebaris $n$ meja, setiap meja menawarkan sisi atas
dan sisi bawahnya ($2n$ tempat duduk), dan kedua meja di ujung
menawarkan satu sisi tambahan masing-masing ($+2$): $2n + 2$ tempat
duduk. Periksa: $4$, $6$, $8$ untuk $n = 1, 2, 3$.

\textbf{9.} Satu baris: $2n + 2 \geq 30$ memerlukan $2n \geq 28$, jadi
$14$ meja. Membelah $14$ meja yang sama menjadi dua baris berisi
$7$: setiap barisnya menampung $2 \times 7 + 2 = 16$, bersama-sama
$32$ --- dua lebih banyak daripada satu baris, karena setiap baris
baru membawa dua \emph{ujung} yang baru. (Setiap pembelahan menambah
$2$ tempat duduk; para penyedia jamuan tahu itu.)

\textbf{10.} Huruf L berukuran $n$ punya $n$ titik ke atas dan $n$
titik melintang, dengan titik pojoknya dipakai bersama:
$n + n - 1 = 2n - 1$ titik. Gambar $1$, $2$, $3$: $1$, $3$, $5$ titik
--- yaitu \emph{\omterm{def:g3:numbers:evenodd}{bilangan ganjil}}. Kalau disarangkan satu mengelilingi
yang lain, huruf L itu mengisi persegi titik $n \times n$: itulah
cerita \cref{pb:g8:equations:1}.

\textbf{11.} $3(n + 2) - 5 = 3n + 6 - 5 = 3n + 1$: pernyataannya
menyusut menjadi rumus yang sudah dibuktikan, jadi ia benar untuk
setiap $n$ --- ditegakkan oleh aljabar, bukan oleh contoh.

\textbf{12.} Ujinya: $4 \geq 2$, $9 \geq 3$, $100 \geq 10$: semuanya
lulus. Tetapi $n = 0.5$ memberi $n^2 = 0.25 < 0.5$: pernyataannya
gagal. Ketiga uji tadi tidak membuktikan apa pun tentang \emph{semua}
bilangan (ketiganya sekadar gagal menemukan masalah); satu contoh
penyangkal itulah yang \emph{membuktikan} pernyataan umumnya salah
(\cref{met:g7:literal:test}). Sebuah huruf mewakili setiap bilangan
--- termasuk \omterm{def:g6:decimals:places}{bilangan desimal}.

\textbf{13.} Daerahnya: $2$ titik $\to 2$; $3$ titik $\to 4$;
$4$ titik $\to 8$; $5$ titik $\to 16$. Polanya membisikkan
``setiap titik melipatduakan hitungannya'': $32$ daerah yang
diharapkan untuk $6$ titik.

\textbf{14.} Pembilangan yang teliti memberi $31$ daerah --- bukan
$32$. Dugaan pelipatduaannya mati: ia selamat empat kali pemeriksaan
lalu gagal pada yang kelima. (Contoh penyangkal ini cukup terkenal
sehingga punya nama, ``perangkap daerah lingkaran''; jangan memakai
jalan pintas titik yang merata --- tiga tali busur yang bertemu di
satu titik akan menggabungkan daerah dan merusak hitungannya.)

\textbf{15.} Pembuktiannya memerlukan alasan \emph{struktural}
mengapa menambah satu titik melipatduakan daerahnya --- dan alasan
itu tidak ada: pelipatduaannya hanyalah kebetulan pada kasus yang
kecil. Rumus batang korek api aman karena masing-masing dibaca
langsung dari susunannya (setiap persegi jelas memakan $3$ batang,
setiap meja jelas menawarkan $2$ tempat duduk): struktur, bukan
pengujian. Ramalan yang terbukti: $3 \times 2\,026 + 1 = 6\,079$ batang.

\end{solution}
